\documentclass[11pt]{article}
\usepackage[margin=1in]{geometry}
\usepackage{amsmath,amssymb,amsthm,hyperref,enumitem}
\newtheorem{theorem}{Theorem}
\newtheorem{lemma}{Lemma}
\title{The reciprocal areas of the lateral harmonic polygons}
\author{Principia Math In Progress}
\date{21 September 2026}
\begin{document}
\maketitle
\noindent\textbf{MathDB 352833.} We prove Conjecture 1 of Garcia, Reznik,
and Roitman \cite[p.~228]{GRR} for every odd number of sides. The source
uses unit circles for polarity; a common radius $\rho>0$ is included below.
The proof has been reviewed by the authoring agent only.

\section*{Task 1: the complete area invariant}
Let the inner ellipse have semiaxes $a\ge b>0$, and put
\[
 c=\sqrt{a^2-b^2},\quad d=\frac{c}{a+b}\in[0,1),\quad
 R=\frac{\rho^2 a}{b^2},\quad \vartheta=\frac{2\pi}{N}.
\]
Consider the convex Poncelet family between this ellipse and its concentric
homothetic outer ellipse, of semiaxes $a/\cos(\pi/N)$ and
$b/\cos(\pi/N)$. The polar polygons are obtained by taking the poles of
the original sides in the two circles of radius $\rho$, centered at
$F_\pm=(\pm c,0)$. Let their ordinary positive areas be $A_\pm$.

\begin{theorem}\label{main}
For every odd integer $N\ge3$, throughout the family,
\begin{equation}\label{invariant}
 \frac1{A_+}+\frac1{A_-}
 =\frac{4(1-2d^2\cos\vartheta+d^4)}{NR^2(1-d^2)^2\sin\vartheta}
   \frac{1+d^{2N}}{1-d^{2N}}.
\end{equation}
In particular this quantity is independent of the polygon's phase.
\end{theorem}

The affine map $(x,y)\mapsto(x/a,y/b)$ takes the two ellipses to
concentric circles. A convex polygon tangent to the inner unit circle and
inscribed in the outer one is regular: each consecutive central angle is
$2\pi/N$. Thus the side lines in the original plane have the form
\begin{equation}\label{sides}
 \frac{x\cos u_j}{a}+\frac{y\sin u_j}{b}=1,\qquad
 u_j=t+j\vartheta,\quad 0\le j<N.
\end{equation}
This also parametrizes the entire family.

For a line $n\cdot X=1$ and a circle $(F,\rho)$, its pole is
$F+\rho^2n/(1-n\cdot F)$: indeed its polar equation is
$(X-F)\cdot(P-F)=\rho^2$. Applying this to \eqref{sides} gives
\[
 P_{+,j}=F_++\rho^2
 \left(\frac{\cos u_j}{a-c\cos u_j},
       \frac{a\sin u_j}{b(a-c\cos u_j)}\right).
\]
Since $a>c$, all denominators are positive. With
$O_+=(c+\rho^2c/b^2,0)$, direct simplification yields
\begin{equation}\label{mobius}
 \frac{P_{+,j}-O_+}{R}
 =\left(\frac{a\cos u_j-c}{a-c\cos u_j},
         \frac{b\sin u_j}{a-c\cos u_j}\right)
 =\frac{z_j-d}{1-dz_j},\qquad z_j=e^{iu_j},
\end{equation}
where the plane on the right is identified with $\mathbb C$.
Here $c/a=2d/(1+d^2)$ and $b/a=(1-d^2)/(1+d^2)$.
For the other focus, $O_-=(-c-\rho^2c/b^2,0)$ and exactly the same
formula holds with $d$ replaced by $-d$. Crucially, the scale $R$ and
the phase $t$ are the same in both formulas.

\begin{lemma}\label{area}
For $N\ge3$, $-1<q<1$, and $z_j=e^{i(t+j\vartheta)}$, the polygon with
vertices $w_j=(z_j-q)/(1-qz_j)$ has area
\begin{equation}\label{areaformula}
 B_N(q,t)=\frac{N}{2}\frac{(1-q^2)^2\sin\vartheta}
 {1-2q^2\cos\vartheta+q^4}
 \frac{1-q^{2N}}{1-2q^N\cos(Nt)+q^{2N}}.
\end{equation}
\end{lemma}
\begin{proof}
The map $z\mapsto(z-q)/(1-qz)$ preserves the unit circle and its cyclic
orientation; its angular derivative is $(1-q^2)/|1-qz|^2>0$.
The resulting polygon is convex and positively oriented. Hence the
shoelace formula gives $B_N=\frac12\operatorname{Im}\sum_j\overline{w_j}w_{j+1}$.
Write $\omega=e^{i\vartheta}$ and
\[
 C=\frac{(1-q^2)(\omega-1)}{1-q^2\omega}.
\]
On the unit circle, $\overline{w_j}=1/w_j$. The following is a rational
identity in $z$:
\begin{equation}\label{partial}
 \frac{(\omega z-q)(1-qz)}{(1-q\omega z)(z-q)}
 =1+C\left(\frac1{1-q\omega z}+\frac{q}{z-q}\right).
\end{equation}
To sum it, use the geometric series and
$\sum_{j=0}^{N-1}z_j^m=0$ unless $N\mid m$. Absolute convergence is
valid because $|q|<1$. Since multiplication by $\omega$ permutes the $z_j$,
\[
 \sum_j\frac1{1-q\omega z_j}=\frac{N}{1-q^Ne^{iNt}},\qquad
 \sum_j\frac{q}{z_j-q}=\frac{Nq^Ne^{-iNt}}{1-q^Ne^{-iNt}}.
\]
Adding these two expressions gives
\[
 \sum_j\overline{w_j}w_{j+1}
 =N\left(1+C\frac{1-q^{2N}}{1-2q^N\cos(Nt)+q^{2N}}\right).
\]
Finally,
\[
 \operatorname{Im}C=
 \frac{(1-q^2)^2\sin\vartheta}{1-2q^2\cos\vartheta+q^4}.
\]
Taking imaginary parts proves \eqref{areaformula}, including $q=0$.
All its denominators are positive.
\end{proof}

\begin{proof}[Proof of Theorem \ref{main}]
Equation \eqref{mobius} gives $A_\pm=R^2B_N(\pm d,t)$, since translation
does not change area. Set
\[
 K=\frac{NR^2}{2}\frac{(1-d^2)^2\sin\vartheta}
 {1-2d^2\cos\vartheta+d^4}(1-d^{2N})>0.
\]
For odd $N$, Lemma \ref{area} gives
\[
 A_+=\frac{K}{1-2d^N\cos(Nt)+d^{2N}},\qquad
 A_-=\frac{K}{1+2d^N\cos(Nt)+d^{2N}}.
\]
The sum of their reciprocals is $2(1+d^{2N})/K$, which is
\eqref{invariant}. If $a=b$, then $d=0$ and the same formula gives the
regular-polygon value; no limiting argument or distinct-focus assumption
is needed. This proves the full conjecture.
\end{proof}

\noindent\textbf{Scope and verification.}
The source already proves the cases $N=3,5$; only the complete odd-$N$
assertion is claimed here. The two polar circles have the same radius,
as in the source's unit-circle convention. The formula also explains the
source's even-$N$ observation: $A_+=A_-$ and, for $d>0$, they vary with
$t$. Symbolic identities, exact algebraic shoelace checks, and supplemental
high-precision geometric checks are in \texttt{verify.py}. The proof above,
rather than finite verification, establishes all parameters.

\section*{Tasks 2--5: proposed follow-ups}
These are proposed mathematical tasks, not additional completed solves.
\begin{enumerate}[label=\textbf{Task \arabic*.},start=2,leftmargin=*]
\item Classify pairs of centers inside the inner ellipse for which a
positive weighted sum of reciprocal polar areas is phase independent.
\item Extend the area formula and the invariant to Poncelet star polygons,
specifying winding numbers and signed versus geometric areas.
\item Compute the area moments $\frac1{2\pi}\int_0^{2\pi}A_+(t)^k\,dt$
for every integer $k$ and describe their eccentricity asymptotics.
\item Determine which reciprocal-area invariants survive when the outer
ellipse is concentric but not homothetic to the inner ellipse.
\end{enumerate}

\begin{thebibliography}{9}
\bibitem{GRR} R. Alves Garcia, D. Reznik, and P. Roitman,
\emph{New Properties of Harmonic Polygons}, Journal for Geometry and
Graphics \textbf{26}(2) (2022), 217--236.
\href{https://www.heldermann-verlag.de/jgg/jgg26/j26h2alve.pdf}{Publisher PDF}.
Current preprint: \href{https://arxiv.org/abs/2112.02545}{arXiv:2112.02545v3}.
Both complete versions were read; the final's Conjecture 1 is on p.~228.
\end{thebibliography}
\end{document}
